\begin{block}{El Lagrangiano Quiral de Yukawa}
  El acoplamiento de Yukawa con sectores escalares y pseudoescalares
  adopta la forma
  \begin{equation}
    \mathcal{L}_{Y} = -\frac{\lambda_{1}v_{1}}{\sqrt{2}}\overline{\psi}\psi
    - \frac{\lambda_{2}v_{2}}{\sqrt{2}}\overline{\psi}\gamma^{5}\psi.
  \end{equation}

  Mediante los parámetros de masa escalar $M$ y pseudoescalar $\tilde{M}$,
  la densidad lagrangiana se expresa en función de las proyecciones
  quirales como
  \begin{equation}
    \mathcal{L}_{\text{Dirac}} = \overline{\psi}_{L}(M+\tilde{M})\psi_{R}
    + \overline{\psi}_{R}(M-\tilde{M})\psi_{L}.
  \end{equation}
\end{block}
